\usetheme[numbering=none,progressbar=none]{metropolis}
% 본문 글씨는 유지하고 수식만 논문 느낌의 serif 수학 글꼴로 사용
\usefonttheme[onlymath]{serif}

\usepackage{iftex}
\usepackage{amsmath,amssymb,bm}
\usepackage{booktabs}
\usepackage{mathtools}
\usepackage{tikz}
\usepackage{xcolor}
\usepackage{graphicx}
\usepackage{adjustbox}

\ifPDFTeX
  \usepackage[T1]{fontenc}
  \usepackage[utf8]{inputenc}
  \usepackage[scaled=0.95]{helvet}
  \usepackage{courier}
  \renewcommand{\familydefault}{\sfdefault}
\else
  \usepackage{fontspec}
  \usepackage{kotex}
\fi

\usetikzlibrary{positioning,arrows.meta,shapes.geometric}

% ====================
% Tool
% ====================

\newcommand{\rmn}[1]{\mathrm{#1}}
\newcommand{\refcite}[1]{{\color{MICGray}[#1]}}
\newcommand{\subsec}[1]{%
    \par\noindent{\normalsize\bfseries #1}%
}

% itemize helper
\newcommand{\items}[1]{%
    \vspace{-0.5em}
    \begin{itemize}
        #1
    \end{itemize}
}

\newcommand{\al}[1]{%
    \begin{align}
        #1
    \end{align}
    \vspace{-0.8em}
}

% ====================
% Column Layout
% ====================

\newlength{\ToolColumnWidth}
\newlength{\ToolArrowWidth}
\newlength{\ToolFlowColumnSep}

\newcommand{\ColumnTitleFont}{\fontsize{9}{11}\selectfont\bfseries}

\definecolor{ColumnBodyBorder}{RGB}{185,185,185}
\definecolor{ColumnTitleBorder}{RGB}{185,185,185}
\definecolor{ColumnTitleFill}{RGB}{226,226,226}
\definecolor{ColumnShadow}{RGB}{0,0,0}

\newcommand{\ColumnArrowWidth}{0.04\linewidth}
\newcommand{\FlowColumnSepValue}{0.35em}
\newcommand{\ArrowColor}{KAISTNavy}
\newcommand{\ArrowSize}{\normalsize}

\newcommand{\Columstart}[1]{%
    \begingroup
    \setlength{\ToolArrowWidth}{\ColumnArrowWidth}%
    \setlength{\ToolFlowColumnSep}{\FlowColumnSepValue}%
    \setlength{\columnsep}{\ToolFlowColumnSep}%
    \pgfmathsetlength{\ToolColumnWidth}{(\linewidth-(((#1-1)*\ToolArrowWidth)+((2*#1-2)*\columnsep)))/#1}%
    \begin{columns}[c,totalwidth=\linewidth]
}

\newcommand{\ColumnEnd}{%
    \end{columns}%
    \endgroup
}

\newcommand{\ColumnBox}[2]{%
    \begin{tikzpicture}
        \node[
            draw=ColumnTitleBorder,
            fill=ColumnTitleFill,
            rounded corners=3pt,
            line width=0.4pt,
            inner xsep=0.7ex,
            inner ysep=0.7ex,
            outer sep=0pt,
            text width=\dimexpr\linewidth-1.4ex\relax,
            align=center,
            anchor=north west
        ] at (0,0) {%
            \begin{minipage}{\dimexpr\linewidth-1.4ex\relax}
                \centering
                {\color{black}\ColumnTitleFont #1}
            \end{minipage}
        };
    \end{tikzpicture}
    \vspace{0.12em}
    \begin{tikzpicture}
        \node[
            draw=ColumnBodyBorder,
            fill=white,
            rounded corners=3pt,
            line width=0.4pt,
            inner xsep=0.05ex,
            inner ysep=0.9ex,
            outer sep=0pt,
            anchor=north west,
            text width=\dimexpr\linewidth-0.1ex\relax,
            preaction={fill=ColumnShadow,opacity=0.12,transform canvas={xshift=1.4pt,yshift=-1.4pt}}
        ] at (0,0) {%
            \begin{minipage}{\dimexpr\linewidth-0.1ex\relax}
                \raggedright
                \begin{itemize}
                    #2
                \end{itemize}
            \end{minipage}
        };
    \end{tikzpicture}
}

\newcommand{\ColumnGroup}[1]{%
    \begin{column}{\ToolColumnWidth}
        #1
    \end{column}
}

\newcommand{\Column}[2]{%
    \ColumnGroup{%
        \ColumnBox{#1}{#2}%
    }%
}

\newcommand{\ColumnArrow}{%
    \begin{column}[c]{\ToolArrowWidth}
        \vfill
        \centering
        {{\ArrowSize\color{\ArrowColor}\(\rightarrow\)}}
        \vfill
    \end{column}
}


\ifPDFTeX\else
  \setsansfont{TeX Gyre Heros}[
    Scale=MatchLowercase
  ]
  \setmainfont{TeX Gyre Heros}[
    Scale=MatchLowercase
  ]
  \setsanshangulfont{Malgun Gothic}[
    Scale=MatchLowercase
  ]
  \setmainhangulfont{Malgun Gothic}[
    Scale=MatchLowercase
  ]
  \setmonofont{Latin Modern Mono}[
    Scale=MatchLowercase
  ]
\fi
\setlength{\parindent}{0pt}

% ====================
% Common Setup
% ====================

% 공통 색상
\definecolor{MICNavy}{RGB}{22,86,156}
\definecolor{MICDeep}{RGB}{16,63,119}
\definecolor{MICSky}{RGB}{86,170,228}
\definecolor{MICSoft}{RGB}{245,249,253}
\definecolor{MICLine}{RGB}{205,219,235}
\definecolor{MICGray}{RGB}{120,120,120}
\colorlet{KAISTNavy}{MICNavy}
\colorlet{SoftBg}{MICSoft}

% 공통 색상 설정
\setbeamercolor{normal text}{fg=black,bg=white}
\setbeamercolor{alerted text}{fg=MICNavy}
\setbeamercolor{example text}{fg=MICDeep}
\setbeamercolor{palette primary}{fg=white,bg=MICNavy}
\setbeamercolor{frametitle}{fg=white,bg=MICNavy}
\setbeamercolor{title}{fg=MICNavy}
\setbeamercolor{subtitle}{fg=MICGray}
\setbeamercolor{progress bar}{fg=MICNavy,bg=MICLine}
\setbeamercolor{progress bar in head/foot}{fg=MICNavy,bg=MICLine}
\setbeamercolor{section in toc}{fg=MICDeep}
\setbeamercolor{block title}{fg=white,bg=MICNavy}
\setbeamercolor{block body}{fg=black,bg=MICSoft}
\setbeamercolor{itemize item}{fg=black}
\setbeamercolor{itemize subitem}{fg=black}

% 목차에서 section 번호를 자동 표시
\setbeamertemplate{section in toc}[sections numbered]

% 공통 글꼴 크기
% 제목 글자 크기: 첫 숫자는 글자 크기, 두 번째 숫자는 줄간격
\setbeamerfont{title}{series=\bfseries,size=\fontsize{20}{24}\selectfont}
% 부제 글자 크기: 첫 숫자는 글자 크기, 두 번째 숫자는 줄간격
\setbeamerfont{subtitle}{size=\fontsize{10}{15}\selectfont}
\setbeamerfont{author}{size=\large}
\setbeamerfont{date}{size=\small}
\setbeamerfont{frametitle}{series=\bfseries,size=\Large}
\setbeamerfont{footline}{size=\fontsize{6.5}{7.5}\selectfont}
\setbeamerfont{CitationNote}{size=\fontsize{6}{7.2}\selectfont}

% ====================
% Theme Setup
% ====================

\metroset{
  titleformat plain=regular,
  titleformat frame=regular,
  sectionpage=none,
  subsectionpage=none,
  block=fill
}

\setbeamertemplate{navigation symbols}{}

% ====================
% Title Page
% ====================

% 표지 로고 크기 배율: 숫자 하나만 바꾸면 전체 크기가 같이 바뀜
\newcommand{\TitleLogoScale}{0.2}
% 표지 로고 위 간격
\newcommand{\TitleLogoX}{0cm}
\newcommand{\TitleLogoY}{-0.315cm}
% 표지 로고 파일을 감싸는 공통 양식
\newcommand{\TitleGraphicBox}[1]{%
  \begin{tikzpicture}[remember picture,overlay]
    \node[anchor=center] at (\TitleLogoX,\TitleLogoY) {%
      \scalebox{\TitleLogoScale}{\includegraphics{#1}}%
    };
  \end{tikzpicture}%
}
\titlegraphic{\TitleGraphicBox{Wordmark-02.png}}

% 표지 전용 상단 박스 바
\newcommand{\TopBar}{%
  \leavevmode%
  \hbox{%
    % ht, dp 값을 바꾸면 박스 바 두께가 바뀜
    \begin{beamercolorbox}[wd=\paperwidth,ht=5.5ex,dp=1.1ex,leftskip=1.0em,rightskip=1.0em]{palette primary}%
    \end{beamercolorbox}%
  }%
}

\makeatletter
% 표지 슬라이드 출력 매크로
\newcommand{\makeTitlePage}{%
  \begingroup
  \setbeamertemplate{headline}{\TopBar}
  \setbeamertemplate{footline}{}
  \begin{frame}[plain]
    % 표지 전체 블록 위치 조정
    % \vspace*{-0.2em}
    \begin{center}
      \ifx\insertsubtitle\@empty\else
      % 부제와 제목 사이 간격
      \vspace{0.1em}
      {\usebeamerfont{subtitle}\usebeamercolor[fg]{subtitle}\insertsubtitle\par}
      \fi
      {\usebeamerfont{title}\usebeamercolor[fg]{title}\inserttitle\par}
      % 제목 블록과 이름 사이 간격
      \vspace{0.8em}
      % 이름 글자 크기
      {{\fontsize{12}{17}\selectfont\insertauthor\par}}
      % 이름과 추가 정보 사이 간격
      \vspace{0.35em}
      % 소속/직위 글자 크기
      {{\fontsize{10}{12}\selectfont\AuthorInfo\par}}
      % 날짜 위 간격
      \vspace{0.5em}
      {\usebeamerfont{date}\insertdate\par}
      % 날짜와 로고 사이 간격
      \inserttitlegraphic
    \end{center}
  \end{frame}
  \endgroup
}

% ====================
% Contents Page
% ====================

% 목차 슬라이드 출력 매크로
\newcommand{\makeContentsPage}{%
  \begingroup
  % 목차 페이지는 위/아래 바를 숨김
  \setbeamertemplate{headline}{}
  \setbeamertemplate{footline}{}
  \begin{frame}{Contents}
    \tableofcontents
  \end{frame}
  \endgroup
}

% ====================
% End Page
% ====================

% 마지막 Q/A 슬라이드 출력 매크로
\newcommand{\makeEndPage}{%
  \begingroup
  \setbeamertemplate{headline}{}
  \setbeamertemplate{footline}{}
  \begin{frame}[plain]
    \begin{center}
      \vfill
      {\usebeamerfont{title}\usebeamercolor[fg]{title}Q/A\par}
      \vfill
    \end{center}
  \end{frame}
  \endgroup
}

% ====================
% Frame Layout
% ====================

% 현재 프레임 제목 저장용 변수
\newcommand{\CurrentFrameTitle}{}

% 슬라이드 안에서 인용문을 넣을 때 사용하는 명령
% 사용 예시: \Citation{[1] Source: Smith et al., 2024}
\newcommand{\Citation}[1]{%
  \begin{tikzpicture}[remember picture,overlay]
    \node[
      anchor=south west,
      xshift=0.5em,
      yshift=2.6ex,
      inner sep=0pt,
      outer sep=0pt
    ] at (current page.south west) {{\color{MICGray}\usebeamerfont{CitationNote}\itshape \shortstack[l]{#1}}};
  \end{tikzpicture}%
}

% 일반 슬라이드 하단 섹션/페이지 번호 바
\newcommand{\SectionBar}{%
  \leavevmode%
  \hbox{%
    % ht, dp 값을 바꾸면 하단 바 높이가 바뀜
    \begin{beamercolorbox}[wd=\paperwidth,ht=1.9ex,dp=0.9ex,leftskip=1.0em,rightskip=1.0em]{palette primary}%
      \strut\usebeamerfont{footline}\thesection.\ \insertsectionhead%
      \hspace{0.6em}%
      |%
      \hspace{0.6em}%
      \CurrentFrameTitle%
      \hfill%
      \strut\insertframenumber/\inserttotalframenumber%
    \end{beamercolorbox}%
  }%
}

\setbeamertemplate{footline}{%
  \SectionBar
}

% 일반 슬라이드 제목 바
\setbeamertemplate{frametitle}{%
  \nointerlineskip
  % ht, dp, leftskip, rightskip 값을 바꾸면 제목 바 크기와 여백이 바뀜
  \begin{beamercolorbox}[wd=\paperwidth,ht=2.5ex,dp=1.1ex,leftskip=1.0em,rightskip=1.0em]{frametitle}%
    % 현재 프레임 제목을 실제 문자열로 저장해서 하단 바에 그대로 표시
    \xdef\CurrentFrameTitle{\insertframetitle}%
    \usebeamerfont{frametitle}\insertframetitle
  \end{beamercolorbox}%
}
\addtobeamertemplate{frame begin}{\gdef\CurrentFrameTitle{}}{}

\renewcommand\[{\begin{equation}}
      \renewcommand\]{\end{equation}}
\makeatother
